% UET Thesis Pandoc LaTeX Template
% Customized for University of Engineering and Technology, VNU

\clearpage
\phantomsection

\setcounter{chapter}{3}
\chapter[Phương pháp xây dựng đường cong elliptic]{Phương pháp xây dựng
đường cong elliptic}

Chương này trình bày trọng tâm lý thuyết của luận văn: các phương pháp
xây dựng đường cong elliptic có bậc nhúng (embedding degree) xác định,
phục vụ cho các ứng dụng mật mã dựa trên phép ghép cặp (pairing-based
cryptography). Chúng ta sẽ đi từ nền tảng lý thuyết của Phương pháp Nhân
phức (CM) và Đa thức lớp Hilbert, đến thuật toán Cocks-Pinch truyền
thống, và cuối cùng là thuật toán Cocks-Pinch Cải tiến được đề xuất
trong luận văn này, với khả năng sinh tham số trường cơ sở có cấu trúc
NTT-friendly nhằm kháng lại tấn công rây trường số tháp (Tower Number
Field Sieve - TNFS).

\section{Ý tưởng cốt lõi}\label{uxfd-tux1b0ux1edfng-cux1ed1t-luxf5i}

Phương pháp Nhân phức (CM) là công cụ nền tảng để \textbf{xây dựng đường
cong elliptic có số điểm định trước}. Thay vì chọn ngẫu nhiên các hệ số
\(a, b\) rồi đếm điểm, CM đi ngược lại: bắt đầu từ bậc nhóm mong muốn
\(\#E(\mathbb{F}_p) = p + 1 - t\) rồi suy ngược ra phương trình đường
cong.

Ý tưởng xuất phát từ lý thuyết số học phức: mỗi đường cong elliptic
\(E\) sở hữu một \textbf{vành nội cấu} (endomorphism ring)
\(\text{End}(E)\). Với đường cong thông thường (không siêu dị), vành này
đẳng cấu với một \textbf{thứ tự trong trường số}
\(\mathbb{Q}(\sqrt{D})\), trong đó \(D < 0\) là số nguyên âm gọi là
\textbf{biệt thức CM}. Giá trị \(D\) đặc trưng hoàn toàn cho ``loại''
đường cong về mặt số học phức.

\subsection{Điều kiện CM và phương trình
Diophantine}\label{ux111iux1ec1u-kiux1ec7n-cm-vuxe0-phux1b0ux1a1ng-truxecnh-diophantine}

Để tồn tại một đường cong elliptic trên \(\mathbb{F}_p\) có vết
Frobenius \(t\) và biệt thức CM là \(D\), điều kiện cần và đủ là phương
trình Diophantine sau phải có nghiệm nguyên:

\[4p = t^2 + |D| \cdot v^2\]

trong đó \(p\) là đặc số trường, \(t\) là vết Frobenius, \(v\) là một số
nguyên dương, và \(D < 0\). Phương trình này \textbf{ràng buộc trực
tiếp} bộ tham số \((p, t, D)\) với nhau; không phải mọi bộ ba nào cũng
có thể thỏa mãn đồng thời.

Ví dụ: Với \(D = -3\) (biệt thức của secp256k1, BLS12-381), phương trình
trở thành \(4p = t^2 + 3v^2\).

\subsection{\texorpdfstring{Đa thức lớp Hilbert
\(H_D(x)\)}{Đa thức lớp Hilbert H\_D(x)}}\label{ux111a-thux1ee9c-lux1edbp-hilbert-h_dx}

Khi đã có \(D\), bước tiếp theo là tính \textbf{Đa thức lớp Hilbert}
(Hilbert class polynomial) \(H_D(x)\). Đây là đa thức hệ số nguyên được
xây dựng từ lý thuyết số học phức mà nghiệm phức của nó chính là các
\textbf{\(j\)-invariant} của tất cả đường cong elliptic phức
\(\mathbb{C}/\mathcal{O}\) với thứ tự CM là \(\mathcal{O}\).

Bậc của \(H_D(x)\) bằng \textbf{số lớp} (class number) \(h(D)\) --- một
bất biến số học của thứ tự \(\mathbb{Z}[\sqrt{D}]\). Với \(|D|\) nhỏ:

\begin{longtable}[]{@{}lllll@{}}
\toprule\noalign{}
\$ & D & \$ & \(h(D)\) & \(H_D(x)\) \\
\midrule\noalign{}
\endhead
\bottomrule\noalign{}
\endlastfoot
3 & 1 & \(x\) & & \\
4 & 1 & \(x - 1728\) & & \\
7 & 1 & \(x + 3375\) & & \\
11 & 1 & \(x + 32768\) & & \\
23 & 3 & \(x^3 + \cdots\) & & \\
\end{longtable}

Với \(|D| = 3\) và \(|D| = 4\), đa thức này bậc 1, nên
\textbf{\(j\)-invariant được xác định ngay lập tức}: \(j = 0\) và
\(j = 1728\) tương ứng. Đây là lý do tại sao secp256k1 (\(j=0\),
\(D=-3\)) và nhiều đường cong BLS có dạng đơn giản \(y^2 = x^3 + b\).

\subsection{Quy trình CM đầy
đủ}\label{quy-truxecnh-cm-ux111ux1ea7y-ux111ux1ee7}

Cho bộ tham số \((p, t, D)\) thỏa mãn \(4p = t^2 + |D|v^2\):

\begin{enumerate}
\def\labelenumi{\arabic{enumi}.}
\tightlist
\item
  \textbf{Tính \(j\)-invariant}: Tìm nghiệm \(j_0 \in \mathbb{F}_p\) của
  \(H_D(x) \pmod{p}\).
\item
  \textbf{Suy ra hệ số đường cong}: Từ \(j_0\), tính
  \(c = j_0 \cdot (1728 - j_0)^{-1} \pmod{p}\), sau đó:

  \begin{itemize}
  \tightlist
  \item
    \(a = 3c\), \(b = 2c\) (trường hợp tổng quát \(j_0 \neq 0, 1728\))
  \item
    \(a = 0\), \(b\) tùy ý (khi \(j_0 = 0\), tức \(D = -3\))
  \item
    \(a\) tùy ý, \(b = 0\) (khi \(j_0 = 1728\), tức \(D = -4\))
  \end{itemize}
\item
  \textbf{Kiểm tra xoắn (Twist test)}: Đường cong
  \(E: y^2 = x^3 + ax + b\) và đường cong xoắn
  \(E': y^2 = x^3 + a\delta^2 x + b\delta^3\) (với \(\delta\) là phần tử
  không phải thặng dư bậc hai) có tổng số điểm \(\#E + \#E' = 2(p+1)\).
  Ta cần chọn đúng đường cong có \(r \mid \#E\).
\item
  \textbf{Tìm điểm sinh (Generator)}: Nhân ngẫu nhiên một điểm trên
  đường cong với hệ số cofactor \(h = \#E / r\) để tìm điểm sinh \(G\)
  có bậc chính xác là \(r\).
\end{enumerate}

\subsection{\texorpdfstring{Trường hợp \(D = -3\): Xoắn
sextic}{Trường hợp D = -3: Xoắn sextic}}\label{trux1b0ux1eddng-hux1ee3p-d--3-xoux1eafn-sextic}

Với \(D = -3\), tình huống phức tạp hơn vì đường cong \(y^2 = x^3 + b\)
có \textbf{6 xoắn} (sextic twists) thay vì chỉ 2. Sáu vết Frobenius khả
dĩ là:

\[t, \quad -t, \quad \frac{t \pm 3v}{2}, \quad \frac{-t \pm 3v}{2}\]

trong đó \(4p = t^2 + 3v^2\). Để xác định xoắn đúng, ta lần lượt thử các
giá trị \(b\) nhỏ và kiểm tra số điểm thực sự của đường cong
\(y^2 = x^3 + b\) bằng cách tính phép nhân vô hướng \([p+1-t']P\) với
điểm ngẫu nhiên \(P\).

Trong cài đặt của luận văn (tệp \texttt{complex\_multiplication.rs}),
hàm \texttt{find\_twist} thực hiện:

\begin{Shaded}
\begin{Highlighting}[]
\CommentTok{// Sáu bậc nhóm khả dĩ cho D = {-}3 (j = 0)}
\KeywordTok{let}\NormalTok{ candidates }\OperatorTok{=}\NormalTok{ [}
\NormalTok{    p1}\OperatorTok{.}\NormalTok{borrowing\_sub(}\OperatorTok{\&}\NormalTok{config}\OperatorTok{.}\NormalTok{t)}\OperatorTok{.}\DecValTok{0}\OperatorTok{,}    \CommentTok{//  p + 1 {-} t}
\NormalTok{    p1}\OperatorTok{.}\NormalTok{carrying\_add(}\OperatorTok{\&}\NormalTok{config}\OperatorTok{.}\NormalTok{t)}\OperatorTok{.}\DecValTok{0}\OperatorTok{,}     \CommentTok{//  p + 1 + t}
\NormalTok{    p1}\OperatorTok{.}\NormalTok{borrowing\_sub(}\OperatorTok{\&}\NormalTok{t3p)}\OperatorTok{.}\DecValTok{0}\OperatorTok{,}         \CommentTok{//  p + 1 {-} (t+3v)/2}
\NormalTok{    p1}\OperatorTok{.}\NormalTok{carrying\_add(}\OperatorTok{\&}\NormalTok{t3p)}\OperatorTok{.}\DecValTok{0}\OperatorTok{,}          \CommentTok{//  p + 1 + (t+3v)/2}
\NormalTok{    p1}\OperatorTok{.}\NormalTok{borrowing\_sub(}\OperatorTok{\&}\NormalTok{t3m)}\OperatorTok{.}\DecValTok{0}\OperatorTok{,}         \CommentTok{//  p + 1 {-} (t{-}3v)/2}
\NormalTok{    p1}\OperatorTok{.}\NormalTok{carrying\_add(}\OperatorTok{\&}\NormalTok{t3m)}\OperatorTok{.}\DecValTok{0}\OperatorTok{,}          \CommentTok{//  p + 1 + (t{-}3v)/2}
\NormalTok{]}\OperatorTok{;}
\CommentTok{// Tìm bậc nào chia hết cho r, rồi tìm b tương ứng}
\end{Highlighting}
\end{Shaded}

\section{Tham số đầu vào của hệ
thống}\label{tham-sux1ed1-ux111ux1ea7u-vuxe0o-cux1ee7a-hux1ec7-thux1ed1ng}

Trước khi đi vào chi tiết thuật toán, phần này giải thích ý nghĩa và cơ
sở lựa chọn của 7 tham số đầu vào trong cài đặt thực tế
(\texttt{cocks\_pinch.rs}):

\begin{Shaded}
\begin{Highlighting}[]
\KeywordTok{let}\NormalTok{ k                    }\OperatorTok{=} \DecValTok{18u64}\OperatorTok{;}
\KeywordTok{let}\NormalTok{ d                    }\OperatorTok{=} \PreprocessorTok{U1024::}\NormalTok{from(}\DecValTok{3}\NormalTok{)}\OperatorTok{;}
\KeywordTok{let}\NormalTok{ target\_r\_bits        }\OperatorTok{=} \DecValTok{512}\OperatorTok{;}
\KeywordTok{let}\NormalTok{ target\_p\_bits        }\OperatorTok{=} \DecValTok{1024}\OperatorTok{;}
\KeywordTok{let}\NormalTok{ max\_attempts         }\OperatorTok{=} \DecValTok{100\_000u64}\OperatorTok{;}
\KeywordTok{let}\NormalTok{ min\_scalar\_two\_adicity }\OperatorTok{=} \DecValTok{32u32}\OperatorTok{;}
\KeywordTok{let}\NormalTok{ min\_base\_two\_adicity   }\OperatorTok{=} \DecValTok{32u32}\OperatorTok{;}
\end{Highlighting}
\end{Shaded}

\subsection{\texorpdfstring{\texttt{k\ =\ 18} --- Bậc nhúng (Embedding
Degree)}{k = 18 --- Bậc nhúng (Embedding Degree)}}\label{k-18-bux1eadc-nhuxfang-embedding-degree}

Bậc nhúng \(k\) xác định \textbf{trường mở rộng} \(\mathbb{F}_{p^k}\)
nơi phép ghép cặp Weil/Tate được tính toán. Giá trị \(k = 18\) được chọn
vì:

\begin{itemize}
\tightlist
\item
  \textbf{Cân bằng bảo mật:} Với \(r \approx 2^{512}\) và
  \(p \approx 2^{1024}\), kích thước trường mở rộng là
  \(p^k \approx 2^{1024 \times 18} = 2^{18432}\) bit. Theo tiêu chuẩn
  NFS hiện tại, bài toán DLP trên \(\mathbb{F}_{p^{18}}\) ở mức bảo mật
  \(\geq 256\) bit --- tương đương với mức bảo mật của khóa
  \(r \approx 2^{512}\).
\item
  \textbf{Hỗ trợ Đa thức Cyclotomic bậc 18:}
  \(\Phi_{18}(T) = T^6 - T^3 + 1\) cho phép sinh \(r\) hiệu quả, có cấu
  trúc NTT-friendly tự nhiên khi \(T \equiv 0 \pmod{2^k}\) (xem §3.3.1).
\item
  \textbf{Thực tế:} \(k = 18\) là bậc nhúng phổ biến trong họ đường cong
  KSS18, phù hợp với mức bảo mật hậu lượng tử cao hơn các đường cong
  BLS12 (dùng trong Ethereum 2.0).
\end{itemize}

\subsection{\texorpdfstring{\texttt{d\ =\ 3} --- Biệt thức CM
(\(|D|\))}{d = 3 --- Biệt thức CM (\textbar D\textbar)}}\label{d-3-biux1ec7t-thux1ee9c-cm-d}

Biệt thức CM \(D = -3\) được chọn vì những ưu điểm nổi bật:

\begin{itemize}
\tightlist
\item
  \textbf{Đa thức Hilbert bậc 1:} \(H_{-3}(x) = x\), nên \(j\)-invariant
  bằng 0 và phương trình đường cong có dạng cực kỳ đơn giản
  \(y^2 = x^3 + b\). Không cần giải đa thức bậc cao.
\item
  \textbf{Phổ biến trong thực tế:} secp256k1 (Bitcoin), BLS12-381
  (Ethereum), và nhiều đường cong quan trọng khác đều dùng \(D = -3\).
\item
  \textbf{Điều kiện CM đơn giản:} Phương trình \(4p = t^2 + 3v^2\) dễ
  giải hơn với hầu hết các giá trị \((t, v)\) nguyên, tăng xác suất tìm
  được \(p\) nguyên tố trong lưới nâng.
\end{itemize}

\subsection{\texorpdfstring{\texttt{target\_r\_bits\ =\ 512} --- Kích
thước bậc nhóm
\(r\)}{target\_r\_bits = 512 --- Kích thước bậc nhóm r}}\label{target_r_bits-512-kuxedch-thux1b0ux1edbc-bux1eadc-nhuxf3m-r}

Tham số này xác định \textbf{mức bảo mật} của bài toán logarit rời rạc
trên đường cong (ECDLP):

\begin{itemize}
\tightlist
\item
  Với \(r \approx 2^{512}\), thuật toán Pollard-rho tốt nhất cần
  \(\approx 2^{256}\) phép tính, đạt mức bảo mật \textbf{256-bit} ---
  vượt tiêu chuẩn NIST SP 800-57 cho ứng dụng đến năm 2030+ và phù hợp
  với kịch bản hậu lượng tử (post-quantum).
\item
  Đây cũng là kích thước của \textbf{khóa riêng} (private key scalar) và
  \textbf{kích thước chữ ký} trong các giao thức Schnorr và ECDSA trên
  đường cong này.
\end{itemize}

\subsection{\texorpdfstring{\texttt{target\_p\_bits\ =\ 1024} --- Kích
thước trường cơ sở
\(p\)}{target\_p\_bits = 1024 --- Kích thước trường cơ sở p}}\label{target_p_bits-1024-kuxedch-thux1b0ux1edbc-trux1b0ux1eddng-cux1a1-sux1edf-p}

Kích thước trường cơ sở tác động đến hai yếu tố:

\begin{itemize}
\tightlist
\item
  \textbf{Bảo mật DLP trên \(\mathbb{F}_{p^k}\):} Với
  \(p \approx 2^{1024}\) và \(k = 18\), kích thước trường mở rộng là
  \(2^{18432}\) bit. Theo ước lượng NFS/TNFS, điều này cho mức bảo mật
  \(\geq 250\) bit --- phù hợp với ứng dụng yêu cầu bảo mật dài hạn.
\item
  \textbf{Hiệu năng:} Kích thước 1024 bit làm cho mỗi phép nhân trường
  tốn gấp đôi so với 512 bit về mặt thời gian CPU, nhưng vẫn trong giới
  hạn thực tế với phép nhân Montgomery đa độ chính xác.
\item
  \textbf{Tỉ số \(\rho\):}
  \(\rho = \log_2 p / \log_2 r = 1024/512 = 2\). Đây là tỉ số điển hình
  của Cocks-Pinch (so với \(\rho = 1.5\) của BLS12 hay \(\rho = 1\) của
  BN).
\end{itemize}

\subsection{\texorpdfstring{\texttt{max\_attempts\ =\ 100\_000} --- Số
lần thử tối
đa}{max\_attempts = 100\_000 --- Số lần thử tối đa}}\label{max_attempts-100_000-sux1ed1-lux1ea7n-thux1eed-tux1ed1i-ux111a}

Giới hạn vòng lặp tìm kiếm an toàn. Trong thực nghiệm, thuật toán CP cải
tiến thường tìm được kết quả trong \textbf{100--6000 lần thử} (tùy cấu
hình NTT), do đó \texttt{100\_000} là giới hạn rất thoải mái đảm bảo
không bao giờ bị timeout trong điều kiện bình thường. Nếu vượt giới hạn
này (không bao giờ xảy ra trong thực tế), hệ thống báo lỗi thay vì chạy
mãi mãi.

\subsection{\texorpdfstring{\texttt{min\_scalar\_two\_adicity\ =\ 32}
--- Two-adicity tối thiểu của
\(r - 1\)}{min\_scalar\_two\_adicity = 32 --- Two-adicity tối thiểu của r - 1}}\label{min_scalar_two_adicity-32-two-adicity-tux1ed1i-thiux1ec3u-cux1ee7a-r---1}

Tham số này ràng buộc \textbf{cấu trúc NTT-friendly của trường vô hướng
\(\mathbb{F}_r\)}:

\[2^{32} \mid (r - 1) \quad \Longleftrightarrow \quad r = d \cdot 2^{32} + 1 \text{ với } d \text{ lẻ}\]

Ý nghĩa thực tế: trường \(\mathbb{F}_r\) chứa một căn nguyên thủy
\(2^{32}\)-th của đơn vị, cho phép thực hiện NTT trên các đa thức bậc
lên đến \(2^{32} \approx 4 \times 10^9\). Đây là yêu cầu cốt lõi để tạo
chứng minh ZK hiệu quả:

\begin{itemize}
\tightlist
\item
  Trong Groth16/PLONK, phép nhân đa thức \(A(x) \cdot B(x)\) với mạch có
  \(2^{20}\) ràng buộc cần NTT bậc \(\geq 2^{21}\) --- thỏa mãn với
  \(2^{32}\).
\item
  BLS12-381 (chuẩn công nghiệp) có \(r\) two-adicity = 32; hệ thống này
  đạt \(\geq 33\), \textbf{vượt chuẩn BLS12-381}.
\end{itemize}

Giá trị 32 đạt được \textbf{không cần rejection sampling} nhờ ràng buộc
\(T \equiv 0 \pmod{2^{11}}\) (vì
\(v_2(r-1) = 3 \cdot v_2(T) \geq 3 \times 11 = 33 \geq 32\)).

\subsection{\texorpdfstring{\texttt{min\_base\_two\_adicity\ =\ 32} ---
Two-adicity tối thiểu của
\(p - 1\)}{min\_base\_two\_adicity = 32 --- Two-adicity tối thiểu của p - 1}}\label{min_base_two_adicity-32-two-adicity-tux1ed1i-thiux1ec3u-cux1ee7a-p---1}

Tương tự, tham số này ràng buộc \textbf{cấu trúc NTT-friendly của trường
cơ sở \(\mathbb{F}_p\)}:

\[2^{32} \mid (p - 1) \quad \Longleftrightarrow \quad p = d \cdot 2^{32} + 1\]

Đây cũng là yêu cầu của thuật toán Miller trong phép ghép cặp: các biểu
thức trung gian (line evaluations) được tính trong \(\mathbb{F}_{p^k}\)
và có thể được tăng tốc bằng NTT nếu \(p\) có cấu trúc NTT. Ngoài ra:

\begin{itemize}
\tightlist
\item
  \(p = d \cdot 2^{32} + 1\) kháng TNFS tốt hơn \(p\) tùy ý (xem §3.3),
  do cấu trúc \(p-1\) không có nhân tử nhỏ ngẫu nhiên mà thuật toán rây
  có thể khai thác.
\item
  Giá trị này đạt được bằng cách mở rộng lưới nâng \((h_t, h_y)\) trong
  hàm \texttt{try\_lift\_to\_prime}.
\end{itemize}

\subsection{Tổng hợp: Đặc trưng đường cong mục
tiêu}\label{tux1ed5ng-hux1ee3p-ux111ux1eb7c-trux1b0ng-ux111ux1b0ux1eddng-cong-mux1ee5c-tiuxeau}

Từ 7 tham số trên, đường cong được xây dựng có đặc trưng:

\begin{longtable}[]{@{}
  >{\raggedright\arraybackslash}p{(\columnwidth - 4\tabcolsep) * \real{0.3333}}
  >{\raggedright\arraybackslash}p{(\columnwidth - 4\tabcolsep) * \real{0.3333}}
  >{\raggedright\arraybackslash}p{(\columnwidth - 4\tabcolsep) * \real{0.3333}}@{}}
\toprule\noalign{}
\begin{minipage}[b]{\linewidth}\raggedright
Thuộc tính
\end{minipage} & \begin{minipage}[b]{\linewidth}\raggedright
Giá trị
\end{minipage} & \begin{minipage}[b]{\linewidth}\raggedright
Ý nghĩa
\end{minipage} \\
\midrule\noalign{}
\endhead
\bottomrule\noalign{}
\endlastfoot
Phương trình & \(y^2 = x^3 + b\) & \(D=-3\), \(j=0\), dạng tối giản \\
Bậc nhúng & \(k = 18\) & Phép ghép cặp trong \(\mathbb{F}_{p^{18}}\) \\
Kích thước \(r\) & 512 bit & Bảo mật ECDLP 256-bit \\
Kích thước \(p\) & 1024 bit & Bảo mật DLP \(\geq 250\)-bit \\
Two-adicity \(r\) & \(\geq 33\) & NTT/FFT trên \(\mathbb{F}_r\) đến bậc
\(2^{33}\) \\
Two-adicity \(p\) & \(\geq 32\) & NTT/FFT trên \(\mathbb{F}_p\) đến bậc
\(2^{32}\) \\
Dạng \(r\) & \(d \cdot 2^{33} + 1\) & NTT-friendly scalar field \\
Dạng \(p\) & \(d \cdot 2^{34} + 1\) & NTT-friendly base field (kháng
TNFS) \\
\end{longtable}

\section{Thuật toán Cocks-Pinch truyền
thống}\label{thuux1eadt-touxe1n-cocks-pinch-truyux1ec1n-thux1ed1ng}

\subsection{\texorpdfstring{Vấn đề: Xây dựng ngược từ
\(r\)}{Vấn đề: Xây dựng ngược từ r}}\label{vux1ea5n-ux111ux1ec1-xuxe2y-dux1ef1ng-ngux1b0ux1ee3c-tux1eeb-r}

Phương pháp CM truyền thống xuất phát từ \(p\) được chọn trước, đòi hỏi
giải phương trình \(4p = t^2 + |D|v^2\) --- một bài toán biểu diễn số
nguyên bởi dạng toàn phương (quadratic form) không hề dễ trong không
gian lớn.

\textbf{Thuật toán Cocks-Pinch (CP)} đảo chiều bài toán: \textbf{bắt đầu
từ \(r\) rồi xây dựng \(p\)}. Đây là cách tiếp cận tự nhiên hơn khi mục
tiêu là đảm bảo \(r\) có kích thước an toàn (ví dụ: 512 bit) với bậc
nhúng \(k\) cho trước.

\subsection{Điều kiện bậc
nhúng}\label{ux111iux1ec1u-kiux1ec7n-bux1eadc-nhuxfang}

Bậc nhúng \(k\) của đường cong là \textbf{số nguyên dương nhỏ nhất} sao
cho \(r \mid p^k - 1\), hay tương đương \(p \equiv 1 \pmod{r}\) khi
\(k=1\), hoặc tổng quát hơn \(p\) là nghiệm của
\(\Phi_k(x) \equiv 0 \pmod{r}\).

Điều kiện cụ thể: \(r \mid \Phi_k(p)\), trong đó \(\Phi_k\) là đa thức
cyclotomic. Điều này có nghĩa là tồn tại \(u\) nguyên tố cùng nhau với
\(k\) sao cho:

\[p \equiv \rho^i \pmod{r} \quad \text{với } \rho \text{ là nghiệm nguyên thủy của } \Phi_k(x) \equiv 0 \pmod{r}\]

\subsection{\texorpdfstring{Thuật toán CP cho
\(k = 18\)}{Thuật toán CP cho k = 18}}\label{thuux1eadt-touxe1n-cp-cho-k-18}

Với \(k = 18\) và \(D = -3\), bậc nhóm
\(r = \Phi_{18}(T) = T^6 - T^3 + 1\) với \(T\) nguyên tùy ý. Các bước:

\textbf{Bước 1 --- Sinh \(r\) từ \(T\):}

\[r = \Phi_{18}(T) = T^6 - T^3 + 1\]

Kiểm tra tính nguyên tố của \(r\) bằng thuật toán Miller-Rabin.

\textbf{Bước 2 --- Tính \(\sqrt{-D} \pmod{r}\):}

Cần \(\beta = \sqrt{-3} \pmod{r}\). Do \(r \equiv 1 \pmod{3}\) (do cấu
trúc của \(\Phi_{18}\)), phần tử này luôn tồn tại. Dùng thuật toán
Tonelli--Shanks.

\textbf{Bước 3 --- Tính \(t_0, y_0\) modulo \(r\):}

Với mỗi \(i\) thỏa \(\gcd(i, k) = 1\) và \(1 \leq i < k\):

\[t_0 = T^i + 1 \pmod{r}, \quad y_0 = \frac{t_0 - 2}{\beta} \pmod{r}\]

\textbf{Bước 4 --- Nâng (lift) lên số nguyên:}

Tìm \(t, y \in \mathbb{Z}\) với \(t \equiv t_0 \pmod{r}\),
\(y \equiv y_0 \pmod{r}\) sao cho:

\[p = \frac{t^2 + 3y^2}{4}\]

là số nguyên và là số nguyên tố. Để làm điều này, thử các dịch chuyển
nhỏ:

\[t = t_0 + h_t \cdot r, \quad y = y_0 + h_y \cdot r, \quad h_t, h_y \in \{-20, \ldots, 20\}\]

cho đến khi \(\frac{t^2 + 3y^2}{4}\) là số nguyên tố đúng kích thước mục
tiêu.

Trong cài đặt (\texttt{cocks\_pinch.rs}), hàm
\texttt{try\_lift\_to\_prime} thực hiện tìm kiếm này:

\begin{Shaded}
\begin{Highlighting}[]
\ControlFlowTok{for}\NormalTok{ ht }\KeywordTok{in} \OperatorTok{{-}}\NormalTok{half\_range}\OperatorTok{..=}\NormalTok{half\_range }\OperatorTok{\{}
    \ControlFlowTok{for}\NormalTok{ hy }\KeywordTok{in} \OperatorTok{{-}}\NormalTok{half\_range}\OperatorTok{..=}\NormalTok{half\_range }\OperatorTok{\{}
        \KeywordTok{let}\NormalTok{ t }\OperatorTok{=}\NormalTok{ apply\_lift(lp}\OperatorTok{.}\NormalTok{t0}\OperatorTok{,}\NormalTok{ r}\OperatorTok{,}\NormalTok{ ht }\KeywordTok{as} \DataTypeTok{i32}\NormalTok{)}\OperatorTok{;}
        \KeywordTok{let}\NormalTok{ y }\OperatorTok{=}\NormalTok{ apply\_lift(lp}\OperatorTok{.}\NormalTok{y0}\OperatorTok{,}\NormalTok{ r}\OperatorTok{,}\NormalTok{ hy }\KeywordTok{as} \DataTypeTok{i32}\NormalTok{)}\OperatorTok{;}
        \CommentTok{// p = (t² + 3y²) / 4}
        \KeywordTok{let}\NormalTok{ numerator }\OperatorTok{=}\NormalTok{ t² }\OperatorTok{+}\NormalTok{ 3y²}\OperatorTok{;}
        \ControlFlowTok{if}\NormalTok{ numerator }\OperatorTok{\%} \DecValTok{4} \OperatorTok{!=} \DecValTok{0} \OperatorTok{\{} \ControlFlowTok{continue}\OperatorTok{;} \OperatorTok{\}}
        \KeywordTok{let}\NormalTok{ p }\OperatorTok{=}\NormalTok{ numerator }\OperatorTok{/} \DecValTok{4}\OperatorTok{;}
        \ControlFlowTok{if}\NormalTok{ is\_prime(}\OperatorTok{\&}\NormalTok{p) }\OperatorTok{\{} \ControlFlowTok{return} \ConstantTok{Some}\NormalTok{(CurveParams }\OperatorTok{\{}\NormalTok{ p}\OperatorTok{,}\NormalTok{ r}\OperatorTok{,}\NormalTok{ t}\OperatorTok{,}\NormalTok{ y}\OperatorTok{,} \OperatorTok{..} \OperatorTok{\}}\NormalTok{)}\OperatorTok{;} \OperatorTok{\}}
    \OperatorTok{\}}
\OperatorTok{\}}
\end{Highlighting}
\end{Shaded}

\subsection{Phân tích độ phức tạp trung
bình}\label{phuxe2n-tuxedch-ux111ux1ed9-phux1ee9c-tux1ea1p-trung-buxecnh}

Mật độ số nguyên tố lân cận \(N\) là khoảng \(1/\ln N\), tức
\(1/1024\ln 2 \approx 1/709\). Với lưới \((h_t, h_y) \in [-20, 20]^2\)
tạo ra \(41 \times 41 = 1681\) ứng viên \(p\) mỗi lần, tỉ lệ thành công
mỗi lần sinh \(r\) là khoảng \(1681/709 \approx 2.4\) ứng viên \(p\)
nguyên tố. Trong thực tế, thường cần khoảng \textbf{100--500 lần sinh
\(r\)} để tìm được bộ tham số hợp lệ.

\section{Thuật toán Cocks-Pinch Cải
tiến}\label{thuux1eadt-touxe1n-cocks-pinch-cux1ea3i-tiux1ebfn}

\subsection{Động lực: Tấn công rây trường số tháp
(TNFS)}\label{ux111ux1ed9ng-lux1ef1c-tux1ea5n-cuxf4ng-ruxe2y-trux1b0ux1eddng-sux1ed1-thuxe1p-tnfs}

Bảo mật của phép ghép cặp phụ thuộc vào độ khó của \textbf{bài toán
logarit rời rạc} (DLP) trên trường mở rộng \(\mathbb{F}_{p^k}\). Năm
2016, Barbulescu và Duquesne chỉ ra rằng thuật toán \textbf{Tower Number
Field Sieve (TNFS)} --- một biến thể của NFS --- có thể giảm đáng kể độ
phức tạp tấn công DLP trên \(\mathbb{F}_{p^k}\) khi \(k\) composite và
\(p\) có cấu trúc đặc biệt. Điều này \textbf{phá vỡ ước lượng bảo mật
ban đầu} của các đường cong như BN256 (từ \textasciitilde128 bit xuống
còn \textasciitilde100 bit hiệu quả).

Để kháng lại TNFS, cần đảm bảo trường cơ sở của đường cong có
\textbf{cấu trúc tốt} ở cả hai trường: trường vô hướng \(\mathbb{F}_r\)
(phục vụ tính toán ZK proof) và trường cơ sở \(\mathbb{F}_p\) (phục vụ
phép ghép cặp).

\subsection{Yêu cầu NTT-friendly}\label{yuxeau-cux1ea7u-ntt-friendly}

Một số nguyên tố \(q\) được gọi là \textbf{NTT-friendly} (hoặc
FFT-friendly) nếu \(q - 1\) chia hết cho một lũy thừa đủ lớn của 2, tức
là:

\[2^s \mid (q - 1) \quad \text{với } s \text{ đủ lớn}\]

Số \(s\) được gọi là \textbf{số hạng hai-adicity} (two-adicity) của
\(q\). Ý nghĩa: trường \(\mathbb{F}_q\) chứa một căn nguyên thủy
\(2^s\)-th của đơn vị, cho phép áp dụng \textbf{Biến đổi Lý thuyết số
(NTT)} --- tức FFT trên trường hữu hạn --- để nhân đa thức bậc cao trong
\(O(n \log n)\) thay vì \(O(n^2)\).

Đây là yếu tố \emph{quyết định} hiệu năng trong: - Sinh chứng minh ZK
(Groth16, PLONK): tính \(H(x) = (A(x) \cdot B(x) - C(x)) / Z(x)\) - Các
cam kết đa thức (KZG, FRI) - Thuật toán MSM (Multi-scalar
multiplication) biến thể NTT

\textbf{Tương quan với TNFS:} Một trường \(\mathbb{F}_p\) với
two-adicity cao có \(p-1 = d \cdot 2^s\), nghĩa là
\(p = d \cdot 2^s + 1\). Cấu trúc này khiến cho tháp trường (tower field
extension) \(\mathbb{F}_{p^k}\) có đặc điểm tốt hơn từ góc độ kháng
TNFS: các thuật toán rây cần khai thác cấu trúc phân tích của \(p-1\) để
tìm quan hệ, do đó khi \(p-1\) có cấu trúc \textbf{không phân tích tùy
tiện} mà theo dạng chuẩn \(d \cdot 2^s + 1\), sự khai thác thông tin này
trở nên khó hơn.

\subsection{\texorpdfstring{Cải tiến 1: NTT-friendly \(r\) qua ràng buộc
\(T\)}{Cải tiến 1: NTT-friendly r qua ràng buộc T}}\label{cux1ea3i-tiux1ebfn-1-ntt-friendly-r-qua-ruxe0ng-buux1ed9c-t}

\textbf{Nhận xét then chốt:} Vì \(r = \Phi_{18}(T) = T^6 - T^3 + 1\), ta
có:

\[r - 1 = T^3(T^3 - 1)\]

Nếu \(T \equiv 0 \pmod{2^k}\), thì \(T^3 \equiv 0 \pmod{2^{3k}}\), và do
đó:

\[v_2(r-1) = v_2(T^3(T^3 - 1)) = 3k\]

(vì \(T^3 \equiv 0 \pmod{2^{3k}}\) và \(T^3 - 1 \equiv -1 \pmod{2}\) là
lẻ).

\textbf{Kết luận:} Để đảm bảo \(\text{two-adicity}(r-1) \geq s\), chỉ
cần \textbf{chọn \(T \equiv 0 \pmod{2^{\lceil s/3 \rceil}}\)}. Đây là
ràng buộc \emph{không cần rejection sampling} --- mọi \(T\) thỏa mãn
điều kiện này đều cho \(r\) với two-adicity ít nhất \(s\).

Trong cài đặt:

\begin{Shaded}
\begin{Highlighting}[]
\CommentTok{// Chọn T là bội số của 2\^{}ceil(s/3): đảm bảo two\_adicity(r{-}1) \textgreater{}= s}
\KeywordTok{let}\NormalTok{ t\_align }\OperatorTok{=}\NormalTok{ min\_scalar\_two\_adicity}\OperatorTok{.}\NormalTok{div\_ceil(}\DecValTok{3}\NormalTok{)}\OperatorTok{;}
\KeywordTok{let}\NormalTok{ step }\OperatorTok{=} \PreprocessorTok{U1024::}\NormalTok{ONE}\OperatorTok{.}\NormalTok{shl(t\_align }\KeywordTok{as} \DataTypeTok{usize}\NormalTok{)}\OperatorTok{;}
\CommentTok{// t\_val luôn là bội số của step}
\KeywordTok{let}\NormalTok{ t\_val }\OperatorTok{=}\NormalTok{ t\_base }\OperatorTok{+} \PreprocessorTok{U1024::}\NormalTok{rand(}\OperatorTok{\&}\NormalTok{t\_steps) }\OperatorTok{*}\NormalTok{ step}\OperatorTok{;}
\KeywordTok{let}\NormalTok{ r }\OperatorTok{=}\NormalTok{ cyclotomic\_phi18(}\OperatorTok{\&}\NormalTok{t\_val)}\OperatorTok{;}
\CommentTok{// Đảm bảo r thỏa mãn: two\_adicity(r{-}1) = 3 * v\_2(T) \textgreater{}= s}
\end{Highlighting}
\end{Shaded}

Với \texttt{min\_scalar\_two\_adicity\ =\ 32}, ta cần
\(\lceil 32/3 \rceil = 11\), tức \(T \equiv 0 \pmod{2^{11}}\). Điều này
\textbf{không làm giảm không gian tham số} đáng kể --- vẫn còn
\(2^{85}/2^{11} = 2^{74}\) ứng viên \(T\) trong dải hợp lệ.

\subsection{\texorpdfstring{Cải tiến 2: NTT-friendly \(p\) qua mở rộng
lưới
nâng}{Cải tiến 2: NTT-friendly p qua mở rộng lưới nâng}}\label{cux1ea3i-tiux1ebfn-2-ntt-friendly-p-qua-mux1edf-rux1ed9ng-lux1b0ux1edbi-nuxe2ng}

Sau khi có \(t_0, y_0 \pmod{r}\), tham số \(p\) được xác định bởi:

\[p = \frac{(t_0 + h_t r)^2 + 3(y_0 + h_y r)^2}{4}\]

Vì \(r \equiv 0 \pmod{2^{3k}}\) (do \(T \equiv 0 \pmod{2^k}\)), các bit
thấp hơn \(3k\) của \(p-1\) chỉ phụ thuộc vào \(t_0, y_0\) (cố định cho
mỗi \(r\)). Tuy nhiên, các bit từ vị trí \(3k\) trở lên bị ảnh hưởng bởi
\(h_t, h_y\).

Để đạt \(\text{two-adicity}(p-1) \geq s_p > 3k\), ta cần mở rộng lưới
\((h_t, h_y)\) ra ngoài phạm vi \(\pm 20\) mặc định:

\[\text{extra} = \max(0,\, s_p - 3k), \quad \text{half\_range} = 20 + 2^{\text{extra}}\]

Điều này đảm bảo có đủ ứng viên \(h_t, h_y\) để \textbf{bao phủ mọi tổ
hợp bit từ vị trí \(3k\) đến \(s_p\)} trong \(p-1\).

\begin{Shaded}
\begin{Highlighting}[]
\KeywordTok{struct}\NormalTok{ LiftParams}\OperatorTok{\textless{}}\OtherTok{\textquotesingle{}a}\OperatorTok{\textgreater{}} \OperatorTok{\{}
\NormalTok{    t0}\OperatorTok{:} \OperatorTok{\&}\OtherTok{\textquotesingle{}a}\NormalTok{ U1024}\OperatorTok{,}
\NormalTok{    y0}\OperatorTok{:} \OperatorTok{\&}\OtherTok{\textquotesingle{}a}\NormalTok{ U1024}\OperatorTok{,}
\NormalTok{    target\_p\_bits}\OperatorTok{:} \DataTypeTok{usize}\OperatorTok{,}
\NormalTok{    min\_base\_two\_adicity}\OperatorTok{:} \DataTypeTok{u32}\OperatorTok{,}  \CommentTok{// s\_p}
\NormalTok{    r\_two\_adicity}\OperatorTok{:} \DataTypeTok{u32}\OperatorTok{,}         \CommentTok{// 3k}
\OperatorTok{\}}

\CommentTok{// Mở rộng lưới tìm kiếm tỉ lệ với khoảng cách giữa s\_p và 3k}
\KeywordTok{let}\NormalTok{ extra }\OperatorTok{=}\NormalTok{ lp}\OperatorTok{.}\NormalTok{min\_base\_two\_adicity}
    \OperatorTok{.}\NormalTok{saturating\_sub(lp}\OperatorTok{.}\NormalTok{r\_two\_adicity)}
    \OperatorTok{.}\NormalTok{min(}\DecValTok{10}\NormalTok{)}\OperatorTok{;} \CommentTok{// giới hạn để lưới không bùng nổ}
\KeywordTok{let}\NormalTok{ half\_range }\OperatorTok{=} \DecValTok{20i64} \OperatorTok{+}\NormalTok{ (}\DecValTok{1i64} \OperatorTok{\textless{}\textless{}}\NormalTok{ extra)}\OperatorTok{;}
\end{Highlighting}
\end{Shaded}

\subsection{Phân tích hiệu năng của thuật toán cải
tiến}\label{phuxe2n-tuxedch-hiux1ec7u-nux103ng-cux1ee7a-thuux1eadt-touxe1n-cux1ea3i-tiux1ebfn}

Bảng dưới so sánh thuật toán CP gốc và CP Cải tiến với cùng mục tiêu
\((r \approx 2^{512}, p \approx 2^{1024}, k = 18)\):

\begin{longtable}[]{@{}
  >{\raggedright\arraybackslash}p{(\columnwidth - 4\tabcolsep) * \real{0.3333}}
  >{\raggedright\arraybackslash}p{(\columnwidth - 4\tabcolsep) * \real{0.3333}}
  >{\raggedright\arraybackslash}p{(\columnwidth - 4\tabcolsep) * \real{0.3333}}@{}}
\toprule\noalign{}
\begin{minipage}[b]{\linewidth}\raggedright
Tham số
\end{minipage} & \begin{minipage}[b]{\linewidth}\raggedright
CP Truyền thống
\end{minipage} & \begin{minipage}[b]{\linewidth}\raggedright
CP Cải tiến (\(s_r = 32, s_p = 32\))
\end{minipage} \\
\midrule\noalign{}
\endhead
\bottomrule\noalign{}
\endlastfoot
Ràng buộc trên \(T\) & Không có & \(T \equiv 0 \pmod{2^{11}}\) \\
two-adicity của \(r-1\) & Ngẫu nhiên (\textasciitilde1-3) & \(\geq 33\)
(đảm bảo) \\
two-adicity của \(p-1\) & Ngẫu nhiên (\textasciitilde1) & \(\geq 32\)
(với lưới mở rộng) \\
Lưới \((h_t, h_y)\) & \([-20, 20]^2\) & \([-21, 21]^2\) (khi
\(s_p \leq 3k\)) \\
Chi phí mỗi lần thử & \textasciitilde50ms & \textasciitilde50ms \\
Số lần thử trung bình & \textasciitilde100--500 &
\textasciitilde3.000--6.000 \\
Thời gian thực tế & \textasciitilde2--10 giây & \textasciitilde30--60
giây \\
\(r\) hỗ trợ NTT? & Không & Có (bậc \(2^{33}\)) \\
\(p\) hỗ trợ NTT? & Không & Có (bậc \(2^{34}\)) \\
\end{longtable}

\emph{Kết quả thực nghiệm điển hình:}

\begin{verbatim}
r two-adicity: 2^33 | (r-1)  →  r = d.2^33 + 1, NTT đến bậc 2^33
p two-adicity: 2^34 | (p-1)  →  p = d.2^34 + 1, NTT đến bậc 2^34
Thời gian tổng: 32.81s
\end{verbatim}

\subsection{So sánh với BLS12-381}\label{so-suxe1nh-vux1edbi-bls12-381}

Đường cong BLS12-381 --- chuẩn công nghiệp hiện tại trong ZK proof ---
có: - \(r\) two-adicity = \textbf{32} (đủ cho ZK circuits đến \(2^{32}\)
ràng buộc) - \(p\) two-adicity = \textbf{32}

Đường cong được xây dựng trong luận văn này đạt \textbf{33--34}, tức là
\textbf{vượt BLS12-381 về cấu trúc NTT} trong khi kích thước khóa lớn
hơn (\(p \approx 2^{1024}\) so với \(p \approx 2^{381}\)), cung cấp mức
bảo mật hậu lượng tử cao hơn trong ngữ cảnh phép ghép cặp.

\subsection{Giới hạn và hướng mở
rộng}\label{giux1edbi-hux1ea1n-vuxe0-hux1b0ux1edbng-mux1edf-rux1ed9ng}

Thuật toán cải tiến vẫn có giới hạn thực tế:

\begin{itemize}
\tightlist
\item
  \textbf{Two-adicity tối đa đạt được}: Khi \(s_p > 3k + 10\), lưới
  \((h_t, h_y)\) phải có \(\text{half\_range} > 2^{10} = 1024\), số lần
  lặp vượt \(10^6\) mỗi ứng viên \(r\), khiến tìm kiếm không thực tế.
  Nếu cần \(s_p\) cực cao (\textgreater50), cần dùng họ đường cong
  chuyên dụng thiết kế từ đầu như Pasta curves (Pallas/Vesta).
\item
  \textbf{Tỉ số \(\rho\)}: Thuật toán CP vốn cho tỉ số
  \(\rho = \log p / \log r \approx 2\), không tối ưu bằng các họ đường
  cong chuyên biệt như BN256 (\(\rho = 1\)). Đây là sự đánh đổi cố hữu
  của phương pháp CP.
\item
  \textbf{Chỉ \(k=18\) và \(D=-3\)} được cài đặt trong hệ thống hiện
  tại. Mở rộng sang \(k = 12, 24\) đòi hỏi thay đổi đa thức cyclotomic
  và công thức CM.
\end{itemize}

\FloatBarrier
